% 实线为请求中的数据，虚线为事后更新；最终宽度约164mm。
\begin{tikzpicture}[
  font=\figurefont\footnotesize,
  rsbox/.style={rectangle,draw=BookRule,fill=BookPaper,
    align=center,text=BookInk,text width=25mm,minimum height=13mm,inner sep=3pt},
  rsflow/.style={-{Stealth[length=1.8mm]},draw=BookMuted,line width=.9pt},
  rslabel/.style={font=\figurefont\footnotesize,text=BookInk,fill=white,inner sep=1pt}
]
  \node[rsbox,fill=FigureInput!12] (input) at (0,0) {请求与目录\\历史、条件、事实};
  \node[rsbox,fill=FigureModel!12] (understanding) at (3.35,0) {理解\\需求与对象依据};
  \node[rsbox,fill=FigureModel!12] (recall) at (6.7,0) {召回\\候选集合};
  \node[rsbox,fill=FigureModel!12] (rank) at (10.05,0) {预估与排序\\候选比较信息};
  \node[rsbox,fill=FigureOutput!10] (display) at (13.4,0) {展示与交互\\实际机会与结果};
  \draw[rsflow] (input.east) -- (understanding.west);
  \draw[rsflow] (understanding.east) -- (recall.west);
  \draw[rsflow] (recall.east) -- (rank.west);
  \draw[rsflow] (rank.east) -- (display.west);
  \node[rsbox,text width=56mm,fill=FigureLoss!10] (state) at (8.35,1.9)
    {累计状态与要求\\预算、曝光、频次及多方目标};
  \draw[rsflow] (state.south) -- (8.35,1) -| (recall.north);
  \draw[rsflow] (8.35,1) -| (rank.north);
  \draw[rsflow] (8.35,1) -| (display.north);
  \node[rsbox,text width=43mm,fill=FigureLoss!10] (learning) at (8.35,-2.15)
    {反馈学习与效果检验\\标签、模型与策略更新};
  \draw[rsflow,rounded corners=1mm]
    (display.south) |- (learning.east);
  \draw[rsflow,dashed,rounded corners=1mm]
    (learning.west) -| (understanding.south);
\end{tikzpicture}
